%% Hand-authored circuitikz figure -- NOT generated by maestro2tex, placed by it via
%% the `order` list in configs/anatop_tb_tia.json.
%% Topology transcribed from the archived netlists under
%%   simarchive/anatop_tb_tia/tia_2026-07-28/Interactive.9/1/<TB>/netlist/input.scs
%% Supplies, the bias trim buses (TCBias/LCBias) and the En pin are omitted; they are
%% identical in every bench and carry no measurement.
\begin{figure}[htbp]
  \centering
  \resizebox{0.50\linewidth}{!}{%
\ctikzset{bipoles/length=1.0cm, resistors/scale=0.85, capacitors/scale=0.85, sources/scale=0.85}
\begin{circuitikz}[
    every node/.append style={font=\footnotesize},
    nlab/.style={font=\scriptsize, inner sep=1.5pt},
    plab/.style={font=\footnotesize\itshape, inner sep=2pt},
    opt/.style={font=\scriptsize, text=black!55, align=center, inner sep=2pt},
    railnode/.style={circle, fill=black, inner sep=1.05pt},
    prb/.style={draw, circle, minimum size=9pt, inner sep=0pt},
  ]

% ================= (a) unity-feedback family =========================
\begin{scope}[shift={(0,0)}]
  \node[op amp] (A) at (0,0) {};
  \node at (A.center) {\footnotesize A2};
  % + input from V_CM
  \draw (A.+) -- (-1.6,-0.3);
  \draw (-1.6,-0.3) to[V, l_=$V_{\mathrm{CM}}$, invert] (-1.6,-2.6) node[ground]{};
  % - input, feedback around the probe
  \draw (A.-) -- (-2.8,0.3);
  \draw (-2.8,0.3) -- (-2.8,2.2);
  \draw (-2.8,2.2) -- (-1.9,2.2);
  \node[prb] (P) at (-1.6,2.2) {};
  \node[nlab, above=2pt] at (P.north) {\texttt{IPRB1}};
  \draw (P.east) -- (1.6,2.2) -- (1.6,0);
  % output
  \draw (A.out) -- (1.6,0);
  \node[railnode] at (1.6,0){};
  \node[nlab, below right] at (1.7,-0.05){\texttt{V\_OutOL}};
  \draw (1.6,0) -- (2.9,0);
  \draw (2.9,0) to[C, l=$C_L$] (2.9,-2.6) node[ground]{};
  \draw (2.9,0) -- (4.4,0);
  \draw (4.4,0) to[I, l=$i_{\mathrm{load}}$] (4.4,-2.6) node[ground]{};
  \node[opt, anchor=north] at (4.4,-3.0) {\texttt{DriveTB} only};
  \node[opt, anchor=north] at (-1.6,-3.0) {DC; a pulse in\\\texttt{StepUnityFeedbackTB}};
  \node[plab, anchor=north west] at (-4.2,-4.3)
        {(a) \texttt{UnityFeedbackTB}, \texttt{DriveTB}, \texttt{StepUnityFeedbackTB}};
\end{scope}

% ================= (b) transimpedance closed loop ====================
\begin{scope}[shift={(0,-7.6)}]
  \node[op amp] (B) at (0,0) {};
  \node at (B.center) {\footnotesize A2};
  % + input from V_CM
  \draw (B.+) -- (-1.6,-0.3);
  \draw (-1.6,-0.3) to[V, l_=$V_{\mathrm{CM}}$, invert] (-1.6,-2.8) node[ground]{};
  % - input is the WE summing node
  \draw (B.-) -- (-3.0,0.3);
  \node[railnode] at (-3.0,0.3){};
  \node[nlab, above=3pt] at (-3.0,0.3){\texttt{WE}};
  \draw (-3.0,0.3) to[I, l_=$i_{\mathrm{we}}$] (-3.0,-2.8) node[ground]{};
  \draw (-3.0,0.3) -- (-4.6,0.3);
  \draw (-4.6,0.3) to[C, l_=$C_{\mathrm{dl}}$] (-4.6,-2.8) node[ground]{};
  % feedback: probe then R_f
  \draw (-3.0,0.3) -- (-3.0,2.2);
  \draw (-3.0,2.2) -- (-2.3,2.2);
  \node[prb] (PB) at (-2.0,2.2) {};
  \node[nlab, above=2pt] at (PB.north) {\texttt{IPRB\_F}};
  \draw (PB.east) -- (-1.2,2.2);
  \draw (-1.2,2.2) to[R, l=$R_f$] (1.6,2.2);
  \draw (1.6,2.2) -- (1.6,0);
  % output
  \draw (B.out) -- (1.6,0);
  \node[railnode] at (1.6,0){};
  \node[nlab, below right] at (1.7,-0.05){\texttt{V\_OUT}};
  \draw (1.6,0) -- (3.0,0);
  \draw (3.0,0) to[C, l=$C_{\mathrm{out}}$] (3.0,-2.8) node[ground]{};
  \node[plab, anchor=north west] at (-4.2,-4.3) {(b) \texttt{TiaFbTB}};
\end{scope}

% ================= (c) open-loop slew ================================
\begin{scope}[shift={(0,-15.0)}]
  \node[op amp] (C) at (0,0) {};
  \node at (C.center) {\footnotesize A2};
  % - input is V_InM, held at V_CM
  \draw (C.-) -- (-4.6,0.3);
  \node[railnode] at (-3.0,0.3){};
  \node[nlab, above=3pt] at (-3.0,0.3){\texttt{V\_InM}};
  \draw (-4.6,0.3) to[V, l_=$V_{\mathrm{CM}}$, invert] (-4.6,-2.6) node[ground]{};
  % differential pulse source from V_InM to V_InP
  \draw (-3.0,0.3) to[V, l_=pulse, invert] (-3.0,-1.5);
  \draw (-3.0,-1.5) -- (-1.6,-1.5) -- (-1.6,-0.3);
  \draw (-1.6,-0.3) -- (C.+);
  % output
  \draw (C.out) -- (1.6,0);
  \node[railnode] at (1.6,0){};
  \node[nlab, below right] at (1.7,-0.05){\texttt{V\_OutOL}};
  \draw (1.6,0) -- (3.0,0);
  \draw (3.0,0) to[C, l=$C_L$] (3.0,-2.6) node[ground]{};
  \node[plab, anchor=north west] at (-4.2,-4.3) {(c) \texttt{SlewOpenLoopTB}};
\end{scope}

\end{circuitikz}
  }
  \caption[{Testbenches used for the transimpedance stage.}]{Testbenches for the transimpedance stage, from their netlists. (a) Unity feedback through a current probe with \(C_L=\SI{5}{\pico\farad}\): \texttt{stb} gives loop gain and phase, adding a swept \(i_{\mathrm{load}}\) makes \texttt{DriveTB}, and a pulsed source \texttt{StepUnityFeedbackTB}. (b) The closed transimpedance loop; the electrode is a bare \SI{1}{\micro\farad} to ground, \emph{not} the Randles cell of Figure~\ref{fig:analog-models} --- the caveat of Section~\ref{sss:tia-loop}. (c) Open-loop differential pulse drive. Supplies, bias trim buses and the enable pin are omitted throughout.}
  \label{fig:tia-benches}
\end{figure}
